% ============================================================================
% Informe Tarea 3 — Sistemas Distribuidos y Paralelos
% Universidad de Concepción, Junio 2026
% ============================================================================
\documentclass[12pt, a4paper]{article}

% ---------- Márgenes ----------
\usepackage[margin=2.5cm]{geometry}

% ---------- Idioma ----------
\usepackage[spanish,es-tabla]{babel}

% ---------- Paquetes obligatorios ----------
\usepackage{graphicx}
\usepackage{booktabs}
\usepackage{amsmath}
\usepackage{amssymb}
\usepackage{hyperref}
\usepackage{float}
\usepackage{pgfplots}
\pgfplotsset{compat=1.18}

% ---------- Configuración de hyperref ----------
\hypersetup{
    colorlinks  = true,
    linkcolor   = blue!70!black,
    citecolor   = green!50!black,
    urlcolor    = blue!60!black
}

% ---------- Espaciado ----------
\usepackage{setspace}
\onehalfspacing
\usepackage{parskip}

\begin{document}

% ============================================================================
%  ENCABEZADO
% ============================================================================
\begin{center}
    {\LARGE \textbf{Tarea 3: Procesamiento de Imágenes en CUDA}} \\[0.5cm]
    {\large Javier Torres \& Juan Umaña \& Benjamin Jimenez} \\[0.2cm]
    {\small Repositorio: \url{https://github.com/khewa11/T3-Paralel}} \\[0.5cm]
    {\large Sistemas Distribuidos y Paralelos} \\
    {\large Universidad de Concepción} \\
    {\large Junio 2026} \\[1cm]
\end{center}

% ============================================================================
%  1. INTRODUCCIÓN
% ============================================================================
\section{Introducción}

El análisis estadístico de grandes colecciones de imágenes constituye un pilar
fundamental en disciplinas como la visión por computador, la compresión con
pérdida y el aprendizaje de representaciones. Una operación recurrente en estos
contextos es el cálculo de la \textbf{matriz de covarianza} a partir de un
conjunto de~$N$ imágenes aplanadas en vectores de dimensión~$D$, donde
$D = W \times H \times C$ representa el producto entre el ancho, el alto y el
número de canales de cada imagen.

\subsection{Formulación matemática}

Sea $\{v^{(k)}\}_{k=1}^{N}$ el conjunto de imágenes aplanadas, con
$v^{(k)} \in \mathbb{R}^{D}$. El procesamiento se descompone en tres etapas:

\begin{enumerate}
    \item \textbf{Vector promedio.} Se calcula el centroide del conjunto:
    \begin{equation}
        \mu_j \;=\; \frac{1}{N} \sum_{k=1}^{N} v_j^{(k)},
        \qquad j = 1, \ldots, D.
        \label{eq:media}
    \end{equation}

    \item \textbf{Matriz centrada.} Cada vector se centra restando el promedio:
    \begin{equation}
        \overline{v}_j^{(k)} \;=\; v_j^{(k)} - \mu_j,
        \qquad k = 1, \ldots, N, \quad j = 1, \ldots, D.
        \label{eq:centrado}
    \end{equation}

    \item \textbf{Matriz de covarianza.} Se obtiene la matriz simétrica
    $\mathbf{C} \in \mathbb{R}^{D \times D}$:
    \begin{equation}
        C_{jj'} \;=\; \frac{1}{N-1} \sum_{k=1}^{N}
            \overline{v}_j^{(k)} \;\overline{v}_{j'}^{(k)},
        \qquad j,j' = 1, \ldots, D.
        \label{eq:covarianza}
    \end{equation}
\end{enumerate}

\subsection{Motivación y limitante física}

El desafío computacional no reside únicamente en la complejidad
$\mathcal{O}(N \cdot D^2)$ del producto exterior acumulado, sino en la
\textbf{limitante física de la memoria de la GPU}. Aplanar miles de imágenes
del dataset DIV2K en una única región contigua de la VRAM excede con creces la
capacidad de tarjetas gráficas de gama media. Por ejemplo, una GPU con
aproximadamente 3.7~GB de VRAM no puede alojar simultáneamente la matriz de
datos completa junto con la matriz de covarianza resultante cuando las
resoluciones se escalan más allá de umbrales modestos.

Este trabajo propone y evalúa una estrategia de \textbf{orquestación asíncrona
mediante CUDA Streams} que particiona los datos en lotes (\textit{batches}),
permitiendo procesar volúmenes de datos que superan la capacidad física de la
VRAM. Se compara esta implementación con una versión tradicional que opera
sobre el \textit{Stream}~0 por defecto, y se analizan las implicaciones del
\textit{latency hiding} sobre el bus PCIe.

% ============================================================================
%  2. DESCRIPCIÓN DE ALGORITMOS
% ============================================================================
\section{Descripción de Algoritmos}

\subsection{Implementación Tradicional (Stream~0)}

La primera implementación sigue el paradigma convencional de programación CUDA:

\begin{itemize}
    \item \textbf{Memoria contigua:} Todas las imágenes se alojan en un único
          buffer de memoria global de la GPU mediante \texttt{cudaMalloc}.
          Igualmente, la matriz de covarianza de tamaño $D \times D$ se aloja
          íntegramente en la VRAM.

    \item \textbf{Stream único (Stream~0):} Todas las transferencias
          \textit{Host}~$\to$~\textit{Device} y las invocaciones de
          \textit{kernels} se serializan en el flujo por defecto, sin
          posibilidad de solapamiento entre cómputo y comunicación.

    \item \textbf{Tiling:} El cálculo de la covarianza emplea \textit{tiling}
          en memoria compartida para maximizar la reutilización de datos dentro
          de cada bloque de hilos, reduciendo los accesos a memoria global.
\end{itemize}

\textbf{Limitación crítica:} Esta implementación falla con un error de
\texttt{Out of Memory} al escalar la resolución de las imágenes (por ejemplo,
a $112 \times 84$ píxeles) debido a que el volumen total de datos excede la
capacidad de la VRAM disponible (3715~MB). La asignación simultánea de la
matriz de datos y la matriz de covarianza resulta inviable bajo estas
condiciones.

\subsection{Orquestación con CUDA Streams}

La segunda implementación resuelve la limitación anterior mediante una
estrategia de particionamiento por lotes con orquestación asíncrona:

\begin{itemize}
    \item \textbf{Particionamiento en lotes (\textit{batches}):} El conjunto
          de~$N$ imágenes se divide en~$S$ lotes de tamaño
          $\lceil N/S \rceil$. Cada lote se procesa de forma independiente,
          acumulando su contribución parcial a la matriz de covarianza.

    \item \textbf{Pinned Memory (\texttt{cudaMallocHost}):} Los buffers del
          lado del host se alojan como memoria \textit{page-locked}, habilitando
          transferencias DMA que no requieren una copia intermedia al
          \textit{staging buffer} del sistema operativo. Esto es requisito
          indispensable para utilizar \texttt{cudaMemcpyAsync}.

    \item \textbf{Transferencias asíncronas (\texttt{cudaMemcpyAsync}):}
          Las copias entre \textit{Host} y \textit{Device} se emiten de forma
          no bloqueante, permitiendo que el motor de copia DMA opere en
          paralelo con la ejecución de \textit{kernels} en otro \textit{stream}.

    \item \textbf{Ocultamiento de latencia (\textit{Latency Hiding}):}
          Al solapar la transferencia del lote~$i+1$ con el cómputo del
          lote~$i$, se busca ocultar la latencia del bus PCIe y reducir el
          tiempo total de ejecución.
\end{itemize}

La acumulación de resultados parciales se realiza directamente sobre la matriz
de covarianza en la GPU, garantizando que el resultado final sea
numéricamente idéntico al de la implementación tradicional.

% ============================================================================
%  3. METODOLOGÍA EXPERIMENTAL
% ============================================================================
\section{Metodología Experimental}

\subsection{Hardware}

La Tabla~\ref{tab:hardware} resume las características del sistema utilizado
para los experimentos.

\begin{table}[H]
    \centering
    \caption{Especificaciones del hardware experimental.}
    \label{tab:hardware}
    \begin{tabular}{@{} l l @{}}
        \toprule
        \textbf{Componente} & \textbf{Especificación} \\
        \midrule
        CPU     & Intel Core i3-1005G1 @ 1.20\,GHz \\
        RAM     & 12\,GB DDR4 \\
        GPU     & NVIDIA GeForce GTX 1650 Ti \\
        VRAM    & 3715\,MB (4\,GB) \\
        Bus     & PCIe 3.0 x16 \\
        SO      & Linux Mint 21.3 \\
        CUDA    & Toolkit 12.x (sm\_75) \\
        \bottomrule
    \end{tabular}
\end{table}

\subsection{Dataset}

Se utilizó el dataset \textbf{DIV2K}, una colección de referencia compuesta por
imágenes de alta resolución ampliamente empleada en tareas de
super-resolución. Para los experimentos, las imágenes fueron redimensionadas y
aplanadas, generando un volumen de datos de aproximadamente 3~MB para el
conjunto de prueba principal.

\subsection{Parámetros del experimento}

El parámetro principal de variación fue el \textbf{número de streams}~$S$,
evaluado en el conjunto $S \in \{1, 2, 4, 8, 16\}$. La configuración con
$S = 1$ representa el caso base en el que la orquestación por lotes no aporta
solapamiento efectivo, funcionando de manera análoga a una ejecución
serializada dentro de la infraestructura de \textit{streams}.

% ============================================================================
%  4. RESULTADOS
% ============================================================================
\section{Resultados}

\subsection{Resultados del Experimento~1}

El Experimento~1 corresponde a la implementación tradicional sobre el Stream~0,
en la cual toda la memoria se aloja de forma monolítica en la VRAM.
La Tabla~\ref{tab:exp1} muestra los tiempos obtenidos para tres resoluciones.

\begin{table}[H]
    \centering
    \caption{Tiempos de ejecución (ms) de la implementación tradicional
             (Stream~0) para distintas resoluciones.}
    \label{tab:exp1}
    \begin{tabular}{@{} l r r r r r @{}}
        \toprule
        \textbf{Resolución} &
        \textbf{$n = W{\times}H{\times}C$} &
        \textbf{H$\to$D (ms)} &
        \textbf{Cómputo (ms)} &
        \textbf{D$\to$H (ms)} &
        \textbf{Total (ms)} \\
        \midrule
        $64 \times 48$   & 9\,216  & 0.74   & 100.79   & 51.41  & 152.93   \\
        $96 \times 72$   & 20\,736 & 1.47   & 1\,007.95 & 540.12 & 1\,549.55 \\
        $112 \times 84$  & 28\,224 & \multicolumn{4}{c}{\textit{Out of Memory}} \\
        \bottomrule
    \end{tabular}
\end{table}

Se observa un crecimiento significativo del tiempo de cómputo y de la copia
\textit{Device}~$\to$~\textit{Host} al pasar de $64 \times 48$ a
$96 \times 72$, consistente con la complejidad cuadrática
$\mathcal{O}(N \cdot D^2)$ del producto exterior acumulado. Al intentar
procesar imágenes de $112 \times 84$ píxeles, la implementación tradicional
falla con un error de \texttt{Out of Memory}, ya que la asignación monolítica
de la matriz de datos y la matriz de covarianza excede la capacidad de la VRAM
(3715~MB).

\subsection{Comparación entre experimentos}

La Tabla~\ref{tab:comparacion} confronta los tiempos totales del Experimento~1
(Stream~0) con los del Experimento~2 (CUDA Streams, $S = 4$) para las tres
resoluciones evaluadas.

\begin{table}[H]
    \centering
    \caption{Comparación de tiempos totales (ms) entre ambos experimentos.}
    \label{tab:comparacion}
    \begin{tabular}{@{} l r r @{}}
        \toprule
        \textbf{Resolución} &
        \textbf{Exp.~1 --- Stream~0 (ms)} &
        \textbf{Exp.~2 --- Streams $S{=}4$ (ms)} \\
        \midrule
        $64 \times 48$   & 152.9   & 1\,997.7 \\
        $96 \times 72$   & 1\,549.6 & 3\,510.3 \\
        $112 \times 84$  & OOM     & 3\,524.8 \\
        \bottomrule
    \end{tabular}
\end{table}

\begin{figure}[H]
    \centering
    \begin{tikzpicture}
        \begin{axis}[
            width=0.85\textwidth,
            height=7cm,
            xlabel={Dimensión $n = W{\times}H{\times}C$},
            ylabel={Tiempo total (ms)},
            xmin=5000, xmax=32000,
            ymin=0, ymax=4200,
            grid=both,
            grid style={line width=0.3pt, draw=gray!30},
            major grid style={line width=0.5pt, draw=gray!50},
            legend pos=north west,
            legend style={font=\small},
            mark size=3pt,
        ]

        % Experimento 1 (Stream 0)
        \addplot[
            thick,
            color=red!80!black,
            mark=square*,
            mark options={fill=red!60},
        ] coordinates {
            (9216,  152.9)
            (20736, 1549.6)
        };
        \addlegendentry{Exp.~1 (Stream~0)}

        % Experimento 2 (Streams S=4)
        \addplot[
            thick,
            color=blue!80!black,
            mark=*,
            mark options={fill=blue!60},
        ] coordinates {
            (9216,  1997.7)
            (20736, 3510.3)
            (28224, 3524.8)
        };
        \addlegendentry{Exp.~2 (Streams $S{=}4$)}

        % Anotación OOM del Exp 1
        \node[anchor=south, font=\footnotesize, text=red!70!black]
            at (axis cs:28224, 200) {\textbf{OOM (Exp.~1)}};

        \end{axis}
    \end{tikzpicture}
    \caption{Tiempo total en función de la dimensión~$n$ para ambos
             experimentos. El Experimento~1 no puede procesar $n = 28\,224$
             (OOM), mientras que la orquestación por lotes del Experimento~2
             completa todas las resoluciones.}
    \label{fig:exp_comparacion}
\end{figure}

A continuación se presentan los resultados detallados del Experimento~2,
correspondiente al procesamiento del dataset de 3~MB con la implementación
basada en CUDA Streams. Se reportan los tiempos de ejecución del \textit{kernel}
en la GPU, el costo de la copia \textit{Device}~$\to$~\textit{Host} y el
tiempo total para cada configuración de~$S$.

\subsection{Tiempos del Experimento~2 variando $S$}

\begin{table}[H]
    \centering
    \caption{Tiempos de ejecución (ms) variando el número de streams~$S$
             sobre el dataset de 3~MB.}
    \label{tab:tiempos}
    \begin{tabular}{@{} r r r r @{}}
        \toprule
        \textbf{Streams ($S$)} &
        \textbf{GPU (ms)} &
        \textbf{Copia D$\to$H (ms)} &
        \textbf{Total (ms)} \\
        \midrule
         1 & 2478.59 & 60.97 & 2539.56 \\
         2 & 2615.94 & 67.06 & 2683.00 \\
         4 & 2365.07 & 55.60 & 2420.67 \\
         8 & 2460.84 & 57.05 & 2517.89 \\
        16 & 2514.05 & 57.88 & 2571.94 \\
        \bottomrule
    \end{tabular}
\end{table}

\subsection{Análisis de Speedup}

Se define el \textit{speedup} relativo tomando como referencia la configuración
con $S = 1$:

\begin{equation}
    \text{Speedup}(S) \;=\; \frac{T_{\text{total}}(S=1)}{T_{\text{total}}(S)}.
    \label{eq:speedup}
\end{equation}

\begin{figure}[H]
    \centering
    \begin{tikzpicture}
        \begin{axis}[
            width=0.85\textwidth,
            height=7cm,
            xlabel={Número de Streams ($S$)},
            ylabel={Speedup},
            xmode=log,
            log basis x=2,
            xtick={1, 2, 4, 8, 16},
            xticklabels={1, 2, 4, 8, 16},
            ymin=0.85,
            ymax=1.15,
            ytick={0.85, 0.90, 0.95, 1.00, 1.05, 1.10, 1.15},
            grid=both,
            grid style={line width=0.3pt, draw=gray!30},
            major grid style={line width=0.5pt, draw=gray!50},
            legend pos=south east,
            legend style={font=\small},
            mark size=3pt,
        ]

        % Línea base (speedup = 1)
        \addplot[
            dashed,
            thick,
            color=black!50,
            domain=0.8:20,
            samples=2,
        ] {1.0};
        \addlegendentry{Línea base ($S_p = 1.0$)}

        % Curva experimental de speedup
        % Speedup: S=1 -> 1.000, S=2 -> 0.9466, S=4 -> 1.0491, S=8 -> 1.0086, S=16 -> 0.9874
        \addplot[
            thick,
            color=blue!80!black,
            mark=*,
            mark options={fill=blue!60},
        ] coordinates {
            (1,  1.0000)
            (2,  0.9466)
            (4,  1.0491)
            (8,  1.0086)
            (16, 0.9874)
        };
        \addlegendentry{Speedup experimental}

        \end{axis}
    \end{tikzpicture}
    \caption{Speedup en función del número de streams. La curva experimental
             oscila en torno a la línea base, evidenciando la ausencia de
             mejora significativa por solapamiento.}
    \label{fig:speedup}
\end{figure}

\subsection{Perfiles de ejecución (Nsight Systems)}

Los perfiles capturados con Nsight Systems para la configuración $S = 4$ sobre
el dataset de 3~MB permiten descomponer el tiempo de ejecución a nivel de
llamadas CUDA y kernels individuales.

\begin{table}[H]
    \centering
    \caption{Distribución del tiempo por llamada CUDA API (Nsight Systems,
             $S = 4$, resolución $96 \times 72$).}
    \label{tab:cuda_api}
    \begin{tabular}{@{} l r r @{}}
        \toprule
        \textbf{Llamada CUDA API} &
        \textbf{Tiempo (ms)} &
        \textbf{\% del total} \\
        \midrule
        \texttt{cudaHostAlloc}                   & 121.93  & 26.5  \\
        \texttt{cudaStreamSynchronize}           & 113.12  & 24.6  \\
        \texttt{accumulateCovKernel<16>}          & 112.81  & 24.5  \\
        \texttt{cudaMemcpy} (total)              &  52.89  & 11.5  \\
        \quad D$\to$H                            &  52.74  & 11.5  \\
        \quad H$\to$D                            &   1.18  &  0.3  \\
        \texttt{cudaMemset}                      &   1.81  &  0.4  \\
        \bottomrule
    \end{tabular}
\end{table}

\begin{table}[H]
    \centering
    \caption{Desglose de kernels ejecutados en la GPU (Nsight Systems,
             $S = 4$, resolución $96 \times 72$).}
    \label{tab:kernels_gpu}
    \begin{tabular}{@{} l r r r @{}}
        \toprule
        \textbf{Kernel} &
        \textbf{Instancias} &
        \textbf{Tiempo (ms)} &
        \textbf{\% GPU} \\
        \midrule
        \texttt{accumulateCovKernel<16>}  & 4 & 112.81  & 99.9  \\
        \texttt{centerBatchKernel}        & 4 &   0.069 &  0.1  \\
        \texttt{accumulateMeanKernel}     & 4 &   0.026 & $<$0.1 \\
        \texttt{finalizeMeanKernel}       & 1 &   0.002 & $<$0.1 \\
        \bottomrule
    \end{tabular}
\end{table}

Del análisis del perfil se extraen tres hallazgos clave:

\begin{enumerate}
    \item El kernel \texttt{accumulateCovKernel<16>} concentra el
          \textbf{99.9\%} del tiempo de GPU, confirmando que el cálculo de la
          covarianza es la operación dominante.
    \item La transferencia \textit{Host}~$\to$~\textit{Device} (1.18~ms)
          es insignificante frente al cómputo, mientras que la copia
          \textit{Device}~$\to$~\textit{Host} (52.74~ms) refleja el costo
          de extraer la matriz de covarianza resultante.
    \item Las llamadas de infraestructura (\texttt{cudaHostAlloc} y
          \texttt{cudaStreamSynchronize}) representan conjuntamente más del
          50\% del tiempo de la API, un overhead significativo que no
          contribuye al cómputo útil.
\end{enumerate}

La razón entre la transferencia de entrada y el cómputo del kernel dominante
resulta:

\begin{equation}
    \frac{T_{\text{H}\to\text{D}}}{T_{\text{cov}}} \;=\;
    \frac{1.18}{112.81} \;\approx\; 1.0\%.
\end{equation}

% ============================================================================
%  5. DISCUSIÓN DE RESULTADOS Y ANÁLISIS CRÍTICO
% ============================================================================
\section{Discusión de Resultados y Análisis Crítico}

Los resultados experimentales revelan un comportamiento aparentemente
paradójico: la orquestación mediante múltiples CUDA Streams no produjo una
reducción significativa del tiempo total de ejecución. Sin embargo, un análisis
riguroso demuestra que este resultado es, de hecho, el comportamiento esperado
dada la naturaleza del cómputo involucrado, y que la verdadera contribución de
la arquitectura por lotes reside en un plano diferente al de la aceleración
temporal.

\subsection{Cuello de botella: dominio del cómputo}

La Tabla~\ref{tab:tiempos} permite cuantificar la distribución del tiempo
entre cómputo y comunicación. Considérese la configuración con $S = 1$: el
tiempo de ejecución en la GPU fue de 2478.59~ms, mientras que la copia
\textit{Device}~$\to$~\textit{Host} requirió únicamente 60.97~ms. Esto
significa que la fase de transferencia representa apenas un
$\mathbf{2.4\%}$ del tiempo total:

\begin{equation}
    \frac{T_{\text{copia}}}{T_{\text{total}}} \;=\;
    \frac{60.97}{2539.56} \;\approx\; 0.024 \;=\; 2.4\%.
\end{equation}

En consecuencia, el cómputo neto domina aproximadamente el \textbf{97.6\%}
del tiempo de ejecución. Esta proporción se mantiene esencialmente constante
a través de todas las configuraciones de~$S$.

El cálculo de la matriz de covarianza es un problema \textbf{estrictamente
\textit{compute-bound}}: la complejidad aritmética del producto exterior
acumulado $\mathcal{O}(N \cdot D^2)$ impone un volumen de operaciones de punto
flotante que satura las unidades de cómputo de la GPU durante la práctica
totalidad del tiempo de ejecución. En este régimen, el solapamiento de
transferencias con cómputo ---la premisa fundamental del \textit{latency
hiding}--- aporta márgenes de mejora \textbf{insignificantes}, ya que el
tiempo ocultable (la transferencia PCIe) es un orden de magnitud inferior al
tiempo irreducible de cómputo.

Esto explica la forma esencialmente \textbf{plana} de la curva de
\textit{speedup} de la Figura~\ref{fig:speedup}: los valores oscilan entre
0.95 y 1.05, fluctuaciones atribuibles a la variabilidad inherente del
\textit{scheduler} de la GPU y del sistema operativo, y no a una tendencia
sistemática de mejora o degradación.

\subsection{Overhead de la infraestructura de streams}

El perfil de Nsight Systems (Tabla~\ref{tab:cuda_api}) revela que las
llamadas de infraestructura asociadas a la orquestación por \textit{streams}
representan una fracción considerable del tiempo total de la API CUDA.
En particular, \texttt{cudaHostAlloc} consume 121.93~ms (26.5\%), ya que la
asignación de memoria \textit{page-locked} implica que el sistema operativo
debe marcar las páginas como no paginables, un proceso costoso que se ejecuta
una sola vez pero que no puede solaparse con cómputo. Por su parte,
\texttt{cudaStreamSynchronize} acumula 113.12~ms (24.6\%), correspondiente a
los puntos de barrera en los que el host espera la finalización de cada lote
en su respectivo \textit{stream}.

Conjuntamente, estas dos llamadas suman más del \textbf{51\%} del tiempo
registrado en la API, constituyendo un overhead de gestión que no aporta
cómputo útil pero que es inherente a la arquitectura de \textit{streams}
asíncronos.

\subsection{Crecimiento cuadrático empírico}

El paso de $64 \times 48$ ($n_1 = 9\,216$) a $96 \times 72$
($n_2 = 20\,736$) en el Experimento~1 permite verificar empíricamente el
crecimiento cuadrático esperado por la complejidad $\mathcal{O}(D^2)$:

\begin{equation}
    \frac{n_2^{\,2}}{n_1^{\,2}} \;=\;
    \frac{20\,736^{\,2}}{9\,216^{\,2}} \;\approx\; 5.06
    \qquad\text{vs.}\qquad
    \frac{T_2}{T_1} \;=\;
    \frac{1\,007.95}{100.79} \;\approx\; 10.0.
\end{equation}

La relación empírica de tiempos ($\approx 10{\times}$) supera la predicción
puramente cuadrática ($\approx 5{\times}$), lo que sugiere la presencia de
factores adicionales como la presión sobre la caché L2, el incremento de la
fragmentación de \textit{tiles} y el crecimiento de la copia
\textit{Device}~$\to$~\textit{Host} proporcional a $D^2$.

\subsection{Mitigación del bus PCIe y la VRAM: éxito arquitectónico}

Si bien el \textit{latency hiding} no logró reducir el tiempo total de
ejecución ---resultado que, como se argumentó, era previsible dada la
naturaleza \textit{compute-bound} del problema---, la orquestación asíncrona
por lotes constituye un \textbf{éxito arquitectónico absoluto} en un aspecto
igualmente crítico: la \textbf{mitigación del límite de la VRAM}.

La implementación tradicional con Stream~0 falla irremediablemente con un
error de \texttt{Out of Memory} al intentar alojar simultáneamente la totalidad
de los datos y la matriz de covarianza en los 3715~MB de VRAM disponibles.
Esta limitación impone un techo inamovible sobre la resolución y el número de
imágenes procesables.

La arquitectura de \textit{streams} con particionamiento en lotes elimina por
completo esta restricción. Al procesar los datos en fragmentos de tamaño
controlado, el consumo pico de VRAM queda acotado por el tamaño de un único
lote más la matriz de covarianza acumulada, independientemente del volumen
total del dataset. En términos prácticos, esto significa que el sistema puede
procesar volúmenes de datos \textbf{arbitrariamente grandes} ---limitados
únicamente por la memoria del host y el tiempo de cómputo--- sin incurrir en
errores de asignación de memoria en la GPU.

La combinación de \texttt{cudaMallocHost} (memoria \textit{pinned}) y
\texttt{cudaMemcpyAsync} garantiza además que las transferencias entre
\textit{Host} y \textit{Device} se realicen de la forma más eficiente posible,
evitando copias intermedias y utilizando directamente el motor DMA del
controlador PCIe.

\subsection{Síntesis del análisis}

El resultado global puede sintetizarse en dos proposiciones complementarias:

\begin{enumerate}
    \item \textbf{El solapamiento de transferencias es irrelevante para
    problemas \textit{compute-bound}.} Cuando el tiempo de cómputo supera en
    uno o más órdenes de magnitud al tiempo de comunicación, el \textit{latency
    hiding} no puede producir aceleraciones perceptibles. La curva de
    \textit{speedup} plana confirma experimentalmente esta predicción teórica.

    \item \textbf{La orquestación por lotes resuelve una limitación
    arquitectónica fundamental.} Independientemente de su impacto en la
    velocidad, el particionamiento permite escalar el procesamiento más allá
    de los límites físicos de la VRAM, transformando un problema intratable
    en uno viable. Este beneficio es ortogonal al \textit{speedup} y
    representa la contribución principal de la implementación con
    \textit{streams}.
\end{enumerate}

% ============================================================================
%  6. CONCLUSIONES
% ============================================================================
\section{Conclusiones}

El presente trabajo implementó y evaluó dos estrategias para el cálculo de la
matriz de covarianza sobre el dataset DIV2K utilizando CUDA: una
implementación tradicional sobre el Stream~0 y una implementación con
orquestación asíncrona mediante múltiples CUDA Streams.

Los resultados experimentales demuestran que:

\begin{itemize}
    \item El cálculo de la matriz de covarianza es un problema
    \textbf{\textit{compute-bound}}, donde el cómputo del producto exterior
    acumulado consume aproximadamente el 97.6\% del tiempo total de ejecución.
    En este régimen, el solapamiento de transferencias PCIe con cómputo no
    produce mejoras significativas en el tiempo de ejecución, como lo evidencia
    una curva de \textit{speedup} esencialmente plana ($S_p \approx 1.0$) para
    $S \in \{1, 2, 4, 8, 16\}$.

    \item La implementación tradicional con asignación monolítica de memoria
    falla por \texttt{Out of Memory} al escalar la resolución de las imágenes,
    imponiendo un límite superior rígido sobre el volumen de datos procesable
    en una GPU con 3715~MB de VRAM.

    \item La orquestación por lotes mediante CUDA Streams, combinada con
    \textit{pinned memory} y transferencias asíncronas, resuelve esta
    limitación de forma definitiva. Al acotar el consumo pico de VRAM al
    tamaño de un único lote, se habilita el procesamiento de datasets
    arbitrariamente grandes sin modificar el hardware subyacente.

    \item Desde una perspectiva de ingeniería de sistemas, la contribución
    principal de la arquitectura de \textit{streams} no es la aceleración
    temporal, sino la \textbf{viabilidad del particionamiento}: convertir un
    problema limitado por la memoria física en uno limitado únicamente por el
    tiempo de cómputo disponible.
\end{itemize}

Estos hallazgos subrayan la importancia de caracterizar correctamente el
\textit{bottleneck} de un algoritmo antes de diseñar estrategias de
optimización. En problemas dominados por el cómputo, las técnicas de
\textit{latency hiding} sobre el bus PCIe son ineficaces para reducir el
tiempo total; no obstante, la infraestructura de \textit{streams} y
transferencias asíncronas sigue siendo indispensable como mecanismo de gestión
de memoria, habilitando la escalabilidad del sistema más allá de las
restricciones físicas del acelerador.

\end{document}
